\Cref{tab:8-threats} 列出主要解释可能失效的原因以及对应对照。测量定义仍见第 3 章；重复构建和精度检查见 \cref{sec:7-8}。

\par\medskip\noindent\begin{minipage}{\linewidth}
\centering\captionsetup{hypcap=false}
\captionof{table}[有效性威胁]{有效性威胁与剩余检查。}
\label{tab:8-threats}
\footnotesize\setlength{\tabcolsep}{4pt}
\begin{tabularx}{\linewidth}{@{}*{3}{>{\raggedright\arraybackslash}X}@{}}
\toprule
受威胁主张 & 替代解释 & 对照与剩余测试 \\
\midrule
执行收益 & 分数、并列或索引状态改变造成表观收益 & CPU 重排序与表分数检查；覆盖有限，进一步测试需重新生成候选数组 \\
预算质量 & 校准样本遗漏困难查询，或工作负载改变 & 独立的带标签测试查询、已记录决策曲线和全池枚举；仍有未测查询子集与未来分布 \\
竞争优势 & 参数网格、慢调用或接口开销改变排序 & 两种召回估计器、保留慢调用、int8 与精度检查、小 batch 比较；其他设备未测 \\
容量／构建主张 & 短生命周期分配逃过峰值采样 & 分配生命周期、重复冷建和 20\,ms 显存采样；更短峰值及完整重训练前沿仍未测 \\
\bottomrule
\end{tabularx}
\end{minipage}\par\medskip

\cref{sec:7-cpu} 中较早的 CPU 比较给出了历史实现在相同召回率下的吞吐差异。独立训练、OpenAI 数据集的列表数不等以及 CPU 计时重复较少，限制了这些比值的解释。针对本文评测 GPU 版本、匹配索引和计时条件的比较仍未完成。

六个数据集不代表六种独立用户工作负载。两个 OpenAI 导出共享 DBpedia 语料，Stella 的查询从语料中留出，而非原始查询集（\cref{sec:3-2}）。比较包含传输与主机接口；内存类型及 C++ 与 Python 调用成本的差异可能影响小 batch 的小幅差距，因此不能全部归因于 GPU 搜索内核（\cref{sec:3-3}）。证据覆盖一张 GPU、静态索引和平方 L2 top-10 搜索。可能仍有近重复，并且缺少相关性标签；更新、并发服务、漂移、十亿规模容量和多 GPU 均未评估。
